\documentclass[11pt]{article}

\usepackage[margin=1in]{geometry}
\usepackage{amsmath,amssymb,booktabs,microtype}
\usepackage[colorlinks=true,linkcolor=blue!50!black]{hyperref}

\title{Discussion: Identifiability and Optimization in Diagonal Kalman Delta Networks}
\author{}
\date{July 29, 2026}

\newcommand{\vct}[1]{\boldsymbol{#1}}
\newcommand{\diag}{\operatorname{diag}}

\begin{document}
\maketitle

\section{Setup}

The diagonal Kalman Delta Network (KDN) predicts a diagonal covariance and uses it to form
an uncertainty-dependent delta-rule write. For key $\vct{k}_t$, transition
$\vct{\alpha}_t$, process noise $\vct{\omega}_t$, and observation noise $r_t$,
\begin{align}
  \widehat{\vct{p}}_t
    &= \vct{\alpha}_t^2\odot\vct{p}_{t-1}+\vct{\omega}_t, \\
  \vct{\kappa}_t
    &= \frac{\widehat{\vct{p}}_t\odot\vct{k}_t}
      {r_t+\sum_i\widehat p_{t,i}k_{t,i}^2}.
\end{align}
The memory update remains a delta update,
\begin{equation}
  \vct{S}_t
  = \vct{D}_t\vct{S}_{t-1}
  + \vct{\kappa}_t
    \left(\vct{v}_t-(\vct{D}_t\vct{S}_{t-1})^\top\vct{k}_t\right)^\top,
  \qquad \vct{D}_t=\diag(\vct{\alpha}_t).
\end{equation}
The implementation additionally tracks diagonal information $\vct{c}_t=1/\vct{p}_t$ with
the calibrated projected update
\begin{equation}
  \vct{c}_t
  = \frac{1}{\widehat{\vct{p}}_t}
  + \frac{d_k}{r_t}\vct{k}_t^2.
\end{equation}
This is a scan-compatible diagonal-information approximation, not an exact diagonal Kalman
posterior.

\section{The Scale Symmetry}

For any constant $\lambda>0$, consider
\begin{equation}
  \vct{p}_t' = \lambda\vct{p}_t,
  \qquad
  \vct{\omega}_t'=\lambda\vct{\omega}_t,
  \qquad
  r_t'=\lambda r_t,
  \qquad
  \mu'=\frac{\mu}{\lambda},
\end{equation}
where $\mu=\vct{c}_0$ is the initial information and hence
$\vct{p}_0=1/\mu$. The predicted covariance scales as
\begin{equation}
  \widehat{\vct{p}}_t'
  = \vct{\alpha}_t^2\odot(\lambda\vct{p}_{t-1})
    +\lambda\vct{\omega}_t
  = \lambda\widehat{\vct{p}}_t.
\end{equation}
Consequently, the gain is invariant:
\begin{equation}
  \vct{\kappa}_t'
  = \frac{\lambda\widehat{\vct{p}}_t\odot\vct{k}_t}
    {\lambda r_t+\sum_i\lambda\widehat p_{t,i}k_{t,i}^2}
  = \vct{\kappa}_t.
\end{equation}
The information update also transforms consistently as
$\vct{c}_t'=\vct{c}_t/\lambda$. Therefore $\vct{\omega}$, $r$, and $\mu$
contain a common covariance-scale degree of freedom that cannot be identified from the memory
output. Positive floors weaken the symmetry near their boundaries but do not remove the general
conditioning problem.

The transition $\vct{\alpha}_t$ is different. It affects both covariance prediction and the
memory transition $\vct{D}_t\vct{S}_{t-1}$, so it is not part of the exact scale symmetry.
Nevertheless, $\vct{\alpha}_t$ and $\vct{\omega}_t$ are partially confounded in the covariance
path because multiple pairs can produce similar $\widehat{\vct{p}}_t$.

\section{What the Noisy-MQAR Results Establish}

The original diagonal model obtains $76.2\pm22.3$ percent over the $5\times5$ initialization
and data grid. Its initialization component of variation (19.2 points) is larger than its data
component (11.4 points). This is not uniform under-capacity: many cells reach 88--96 percent,
while a few fail severely.

The most informative intervention is isotropic gain initialization. Zeroing the output map of
the process-noise projection makes $\vct{\omega}_t$ uniform across channels at initialization,
after which anisotropy is learned. At the same learning rate and with the same optimizer, this
changes performance from
\begin{equation}
  76.2\pm22.3 \quad\longrightarrow\quad 83.7\pm10.4,
\end{equation}
and halves the initialization variation from 19.2 to 9.0 points. It also rescues both severe
collapse cells. Thus the current evidence identifies random initial gain anisotropy as the
primary failure, not the nominal learning rate or scale symmetry.

The scale symmetry remains relevant as a secondary optimization defect. It creates a flat
direction, gives equivalent physical models different raw parameter values, and lets coupled
quantities evolve under different optimizer geometries. In the current Muon setup, matrix
weights use Muon with adjusted learning rates, while biases and the one-dimensional prior
parameter use AdamW. A single printed learning rate therefore does not imply uniform motion of
the physical uncertainty variables.

The prior $\mu$ is unlikely to explain late failures. It only controls the initial information,
and its influence should decay rapidly over a long sequence when process noise has a positive
floor. Learning it adds a weakly identified parameter with little expected benefit.

\section{Ways to Remove or Reduce the Symmetry}

\subsection{Fix the observation-noise scale}

The strongest and simplest gauge choice is
\begin{equation}
  r_t\equiv1, \qquad \mu\equiv1.
\end{equation}
The model then learns process uncertainty in units of observation noise:
\begin{equation}
  \vct{\omega}_t
  = \omega_{\min}+\operatorname{softplus}(f_\omega(\vct{x}_t)).
\end{equation}
Because $r$ can no longer be rescaled, the common covariance-scale symmetry is removed. This
parameterization is the preferred baseline unless token-dependent observation reliability is
shown empirically to matter.

\subsection{Learn bounded relative observation reliability}

If heteroscedastic observation noise is useful, it should change confidence relative to a fixed
reference rather than learn an unrestricted global scale. One option is
\begin{equation}
  r_t=\exp\!\left(s\tanh f_r(\vct{x}_t)\right),
\end{equation}
with fixed $s$. For $s=\log4$, $r_t\in[1/4,4]$. The projection should have no learnable global
bias, and the prior can remain fixed at $\mu=1$. A soft anchoring penalty can further control
drift,
\begin{equation}
  \mathcal{L}_{\mathrm{scale}}
  = \gamma\left(\mathbb{E}[\log r_t]\right)^2.
\end{equation}
This reduces rather than mathematically eliminates every possible input-dependent rescaling,
but it prevents the practically harmful unbounded common-scale direction.

\subsection{Parameterize a stationary uncertainty level}

The $\vct{\alpha}$--$\vct{\omega}$ interaction can be made more interpretable with
\begin{equation}
  \omega_{t,i}
  = \omega_{\min}+(1-\alpha_{t,i}^2)\rho_{t,i},
  \qquad
  \rho_{t,i}=\operatorname{softplus}(f_\rho(\vct{x}_t)).
\end{equation}
Here $\alpha$ controls retention and $\rho$ is the uncertainty level toward which prediction
moves. This is exactly the completed \texttt{DiagKLA1} experiment. Coupling alone changes
$76.2\pm22.3$ to $76.9\pm20.9$, which is effectively a wash. The completed
\texttt{DiagKLA1-zi} experiment combines coupling with isotropic gain initialization and obtains
$80.1\pm13.3$, below the simpler \texttt{DiagKLA-zi} result of $83.7\pm10.4$. Thus this
parameterization has already been tested: it neither replaces isotropic initialization nor
improves on free additive process noise once that initialization is used.

\section{Recommended Ablation}

The transition-coupled process-noise comparison is already complete:
\begin{center}
\begin{tabular}{lcc}
\toprule
Variant & Free additive $\omega$ & $(1-\alpha^2)\rho$ \\
\midrule
Random anisotropic initialization & $76.2\pm22.3$ & $76.9\pm20.9$ \\
Isotropic gain initialization & $83.7\pm10.4$ & $80.1\pm13.3$ \\
\bottomrule
\end{tabular}
\end{center}
The remaining experiment should isolate covariance-scale choices while keeping free additive
$\omega$ and isotropic gain initialization fixed:
\begin{center}
\begin{tabular}{llll}
\toprule
Variant & $\mu$ & $r_t$ & $\omega_t$ \\
\midrule
A & learned & learned positive & free additive \\
B & fixed at 1 & learned positive & free additive \\
C & fixed at 1 & fixed at 1 & free additive \\
D & fixed at 1 & bounded relative & free additive \\
\bottomrule
\end{tabular}
\end{center}
Run at least the same $5\times5$ initialization--data grid. Report mean, total standard
deviation, initialization variation, non-finite runs, and convergence time. During training,
log $\omega$, $r$, effective write strength, gain anisotropy, the minimum channel gain, and
gradient norms for the process-noise, observation-noise, prior, and transition parameters.

If C matches or improves A, the production model should remove both $r_{\mathrm{proj}}$ and
$\mu_{\mathrm{param}}$. If D improves C, retain only bounded relative observation reliability.
Only after fixing this gauge is a separate uncertainty-parameter learning-rate sweep readily
interpretable.

\section{Conclusion}

The diagonal KDN has a genuine common-scale non-identifiability in
$(\vct{\omega},r,\mu)$, but the completed noisy-MQAR experiments show that random initial
anisotropy is the dominant observed instability. The appropriate design is therefore
twofold: start the gain isotropically, and fix a covariance unit---preferably $r=1$ and
$\mu=1$---so the model learns only uncertainty ratios that affect its predictions.

\end{document}
